% Generated by Listas/make_statements.py from pratica-lecture-08.md, item 3. Do not edit.
\textbf{Source.} Tadelis (2013), \emph{Game Theory: An Introduction}, Exercise 15.3 ``Trust Thy Neighbor'', pp.~315--316 {[}PDF pp.~334--335{]}. Reproduced verbatim. Answer: none in the book.



\begin{quote}

\textbf{Trust Thy Neighbor:} A student, player 1, has to hand in a problem set at the other end of (a large) campus but needs to rush into an exam. She has two options. She can deliver the problem set after the exam (call this option \(L\)) and incur a late penalty. Alternatively she can give the problem set to player 2, a random student who happens to be next to player 1 (call this option \(S\)). Player 2 can then either deliver the problem set on time (call this option \(D\)) or throw it away in the nearest trash can (call this option \(T\)). For player 1 the payoff is 1 if the problem set is delivered on time, \(-1\) if it is thrown away, and 0 if it is delivered late. The payoffs for player 2 are \(x\) if he delivers and \(y\) if he throws it away.



\begin{enumerate}

\def\labelenumi{\alph{enumi}.}

\item

  Draw the game tree that represents this game. Using equilibrium analysis, what conditions on \(x\) and \(y\) would justify player 1 trusting player 2 to deliver the problem set?

\item

  Now assume that some proportion \(p\) of the population of students are ``nice guys'' (\(N\)), for which \(x = 1\) and \(y = 0\), while a proportion \(1 - p\) of the population are ``jerks'' (\(J\)), for which \(x = 0\) and \(y = 1\). Thus we can think that before the game is played, Nature first chooses the type of player 2 that player 1 will meet, according to the probabilities previously described. Draw the game tree.

\item

  Now assume that for (b) \(p = \frac{3}{4}\). What are the pure-strategy Bayesian Nash equilibria of this game? Are there Bayesian Nash equilibria that are not perfect Bayesian equilibria?

\item

  For which values of \(p\) are all the Bayesian Nash equilibria also perfect Bayesian equilibria?

\end{enumerate}

\end{quote}

